\documentclass[main]{subfiles}
\usepackage[sols]{workbook}
\numbering{ws3-2}

\begin{document}

\ifSubfilesClassLoaded{%
  \chaptermark{\chapterthree}
}{}

\section{Slope Fields} \label{ws3-2}

\begin{framed}
\centerline{\textbf{Lesson Goals}}

You should be able to:
\begin{itemize}
\item Sketch a slope field for a given differential equation, and use it to sketch specific solutions and guess the general solution
\item Use a slope field to determine how different solutions of a differential equation relate to each other graphically
\item Use slope fields to determine the long-term behavior of a solution to a differential equation
\item Use Euler's method to approximate a solution at a specific instant $t$ given an initial condition at $t=0$
\end{itemize}
\end{framed}


\begin{enumerate}
    
 \item 
  \note{I usually have students try to guess the general solution of each
  DE after they've sketched the slope field.  It often takes them a few
  tries to get the right answer, so this is a good opportunity to
  reinforce the message that it's easy to check whether a function is
  truly a solution of a DE, simply by plugging it in.}
  
  Draw slope fields for the following differential equations:

  \probonly{\begin{multicols}{2}}
  \begin{enumerate}
    \item 
    \begin{problem}[2.1in]
    $\displaystyle \frac{dy}{dt} = 1$.
    \end{problem}

    \begin{solution}

    \includegraphics[width=2in]{ws/images/draw-slope-fields-1}
    \end{solution}

    \item
    \begin{problem}[\vfill]
    $\displaystyle \frac{dy}{dt} = -t$.
    \end{problem}

    \begin{solution}

    \includegraphics[width=2in]{ws/images/draw-slope-fields-2}
    \end{solution}

    \item
    \begin{problem}[2.1in]
    $\displaystyle \frac{dy}{dt} = -y$.
    \end{problem}

    \begin{solution}

    \includegraphics[width=2in]{ws/images/draw-slope-fields-3}
    \end{solution}        

    \item
    \begin{problem}[\vfill]
    $\displaystyle \frac{dy}{dt} = -\frac{t}{y}$.
    \end{problem}

    \begin{solution}

    \includegraphics[width=2in]{ws/images/draw-slope-fields-4}
    \end{solution}  
  \end{enumerate}
  \probonly{\end{multicols}}
  \pspace{\vfill}
  \pspace{\vfill}

  \probonly{Now, go back and try to figure out the general solution of
  each differential equation. Also try to figure out the specific solution that satisfies the initial condition $y(0)=1$.} 


  \pspace{\newpage}

%%%%%%%%%%%%%%%%%%%%%%%%%%%%%%%%%%%%%%%%%%%%%%%%%%%%%%%%%%%%%%%%%%%%%%%%%%%%%%%%%%%%


  \item 
  \begin{enumerate}
    \item
    \begin{problem}[1.0in]
    If you have a differential equation of the form $\displaystyle \frac{dy}{dt} =
    f(t)$ (that is, the right-hand side depends only on $t$, not on $y$; an
    example is $\displaystyle \frac{dy}{dt} = t^2 + 1$), what can you say about its
    slope field?  How do different solutions of the differential
    equation relate to each other?
    \end{problem}

    \begin{solution}
    When drawing the slope field, all of the slopes along a fixed
    vertical line will have the same slope; after all, the points along
    a fixed vertical line all have the same $t$-value.  For instance,
    look at the slope field for {{\#1}(b)}; all of
    the slopes along a vertical line (such as the vertical line $t = 1$)
    are the same.

    From this, we can conclude that different solutions of the
    differential equation are simply vertical translates of each other.
    \end{solution}
    
    \item
    \begin{problem}[1.0in]
    If you have a differential equation of the form $\displaystyle \frac{dy}{dt} =
    f(y)$ (that is, the right-hand side depends only on $y$, not on $t$),
    what can you say about its slope field?  How do different solutions
    of the differential equation relate to each other?
    \end{problem}

    \begin{solution}
    The idea is very similar to the previous part.  When drawing the
    slope field, all of the slopes along a fixed \textit{horizontal}
    line will have the same slope.  (For instance, look at the slope
    field in {\# 1}(c) above.)   Therefore, we can translate
    a given solution horizontally to get another solution.
    \end{solution}
    
  \end{enumerate}

%%%%%%%%%%%%%%%%%%%%%%%%%%%%%%%%%%%%%%%%%%%%%%%%%%%%%%%%%%%%%%%%%%%%%%%%%%%%%%%%%%%%
  
  \item 
  \note{This problem can be nice for talking about whether different
  solutions of the same differential equation can cross or merge.}
  
  \begin{enumerate}
    \item 
    \begin{problem}[1.85in]
    Draw the slope field for the differential equation $\displaystyle \frac{dy}{dt} =
    1 - y$.  Sketch two solutions to the equation.
    \end{problem}

    \begin{solution}
    The plot on the left shows the slope field.  The plot on the right
    shows two solutions.

    \includegraphics[width=2in]{ws/images/draw-slope-fields-5} \hspace{0.5in} 
    \includegraphics[width=2in]{ws/images/draw-slope-fields-5-sols}
    \end{solution}

    \item

    \note{Although it's fine for them to simply try each one in turn, it's
    nice if they can look at the slope field and figure out which choice
    looks the most plausible.  (This is easy if they see how this slope
    field relates to the one for $\frac{dy}{dt} = -y$.)}
    
    \begin{problem}[.3in]
    Which of the following is a solution to $\displaystyle \frac{dy}{dt} = 1 - y$?

\begin{multicols}{2}
    \begin{enumerate}[label=\protect\maybecircle{\roman*.}]
    \plainitem $y = Ce^{-t}$
    \plainitem $y = Ce^t -t$
    \plainitem $y = Ce^{-t} -1$
    \circleitem $y = Ce^{-t} +1$
    \plainitem $y = Ce^t +1$
  \end{enumerate}
\end{multicols}

  \end{problem}

  \begin{solution}
  You might have noticed that the slope field for $\frac{dy}{dt} = 1 -
  y$ looks exactly like the slope field for $\frac{dy}{dt} = -y$,
  except that it's shifted up by 1.  So, the solution curves should be
  the solution curves of $\frac{dy}{dt} = -y$ 
  (and we know that those solutions are $Ce^{-t}$),
  shifted up by 1.  Thus, a good guess is $y = Ce^{-t} +
  1$, and we can check that this is correct: if $y = C e^{-t} + 1$,
  then $\frac{dy}{dt} = - C e^{-t}$, which is the same as $1 - y = - C
  e^{-t}$.
  \end{solution}
  
  \probonly{What information would you need to solve for the value of $C$?}
\end{enumerate}



%%%%%%%%%%%%%%%%%%%%%%%%%%%%%%%%%%%%%%%%%%%%%%%%%%%%%%%%%%%%%%%%%%%%%%%%%%%%%%%%%%%%

\pspace{\newpage}
\item
\begin{problem}
  Suppose that population density of sharks is modeled by $\displaystyle \frac{dP}{dt} =
  \frac{5}{8} P \left( 1 - \frac{P}{10} \right)$, where $P(t)$ represents density in sharks per square 
  kilometer and $t$ is measured in years.
\end{problem}

\begin{enumerate}
  \item 
  \begin{problem}[2.7in]
    Sketch the slope field and a few representative solutions of the differential equation on the same set of axes.  Identify
    any constant solution.
  \end{problem}
 
  \begin{solution}
  The constant solutions are $y=0$ and $y=10$.  We will talk
  more about these kinds of solutions next class.

%\begin{minipage}{1in}
\begin{tikzpicture}[declare function={f(\x,\y)=(5/8)*\y*(1-\y/10);},scale=0.45]
\def\xmax{20} \def\xmin{-1}
\def\ymax{15} \def\ymin{-1}
\def\nx{15}
\def\ny{15}

\pgfmathsetmacro{\hx}{(\xmax-\xmin)/\nx}
\pgfmathsetmacro{\hy}{(\ymax-\ymin)/\ny}
\foreach \i in {0,...,\nx}
\foreach \j in {0,...,\ny}{
\pgfmathsetmacro{\yprime}{f({\xmin+\i*\hx},{\ymin+\j*\hy})}
\draw[blue,shift={({\xmin+\i*\hx},{\ymin+\j*\hy})}] 
(0,0)--($(0,0)!2mm!(.1,.1*\yprime)$);
}

% solutions
\def\yo{1}
\draw[green,very thick] plot[domain=0:20] (\x,{10});
\draw[magenta, very thick] plot[domain=0:20] (\x,{0});
\draw[blue,very thick] plot[domain=0:20] (\x,{140/(-4*exp(-5*\x/8)+14)});
\draw[red,very thick] plot[domain=0:20] (\x,{20/(8*exp(-5*\x/8)+2)});

\draw[->] (\xmin-.5,0)--(\xmax+.5,0) node[below right] {$x$};
\draw[->] (0,\ymin-.5)--(0,\ymax+.5) node[above left] {$y$};
%\draw (current bounding box.north) node[above]
%{Slope field of \quad $y'=x+y$.};
\end{tikzpicture}
%\end{minipage}


    
  %\includegraphics[width=2.5in]{ws/images/sharks.png}
    
  \end{solution}
    
  \item
  \begin{problem}
    On your qualitative picture above, draw a sketch of the particular solution that 
    satisfies the initial condition $P(0)=2$ (i.e., the initial
    population density right now is $2$ sharks per square km).
  \end{problem}


  \begin{solution}
    The particular solution to the initial value problem in part b) is drawn in red.
  \end{solution}

  
\end{enumerate}
    
The rough sketch you drew in the previous part is enough to see what the long-term behavior of 
the population density is, but not enough to give
you useful \emph{quantitative} information about the population density: for example, what is the population density two years from now?
  
\begin{enumerate}[resume]
    
  \item \label{qualitative-linear-approx-at-2}
  \begin{problem}[1.5in]
    Assume that the current population density is $2$ sharks per square km.  
    Use linear approximation to approximate the population density $4$ years 
    from now---that is, approximate $P(4)$, assuming that $t = 0$ corresponds
    to the present time.  (Hint: you know $P(0)$; what is $P'(0)$?)
  \end{problem}
    
    
  \begin{solution}
    We know that $P(0)=2$. The differential equation tells us that $P' = \frac{5}{8}P(1-P/10)$, 
    which means that:
    $$
      P'(0) = \frac{5}{8} P(0) \left(1 - \frac{P(0)}{10} \right) \; = \; \frac{5}{8} (2) \left( 1 - \frac{2}{10} \right) \; = \; 1.
    $$
      Now we can use linear approximation (centered at $t=0$) to approximate:    
    $P(4) \approx P(0) + P'(0)(4-0) = 2 + (1)(4-0) = 6$     
    Therefore, we expect a population density of about 6 sharks per square km four years from now.
  \end{solution}

\pspace{\newpage}      
  \item
  \begin{problem}[1in]
    Your roommate wants to use the same technique as in \ref{qualitative-linear-approx-at-2} to estimate $P(100)$ by approximating
    \[P(100) \approx P(0) + P'(0)(100-0) \; = \; 2 + (1)(100 - 0) \; = \; 102.\]
    How might you convince your roommate that this is not a reasonable estimate?
  \end{problem}   
  
  \begin{solution}
    If we center our linear approximation at $t = 0$, then $t = 100$ is very ``far" from 
    the center.  Based on everything we saw during the series unit, we shouldn't have much
    confidence in an approximation of $P(100)$ that's based solely on $P(0)$ and $P'(0)$.
    Forecasting how things will look 100 years from now based upon a single data point 
    doesn't feel safe.
 
    More importantly, your roommate's estimate predicts that the population will grow past
    the carrying capacity of 10 sharks per square km.  This goes against what the model 
    predicts---see also the qualitative analysis in Part (a) above.
    
    Linear approximation predicts a future value based on the assumption that the rate of 
    change is constant.  But we know that population density grows more slowly as the density
    draws closer to 10.
  \end{solution}
  
  \item
  \begin{problem}[1.0in]
    Two students are both trying to find an approximate value for $P(4)$.
    \begin{itemize}
      \item Student {\#1} uses the linear approximation from Part (c) above.
      \item Student {\#2} decides to estimate $P(4)$ in stages:  First, use
      linear approximation at $t = 0$ to estimate $P(2)$:
      $$
        P(2) \; \approx \; P(0) + P'(0)(2 - 0) \; = \; 4.
      $$
      Then, the differential equation provides an estimate of $P'(2)$:
      $$
        P'(2) \; = \; \frac{5}{8} P(2) \left[ 1 - \frac{P(2)}{10} \right] \; \approx \; 1.5.
      $$
      Finally, use linear approximation at $t=2$ to estimate
      $$
        P(4) \; \approx \; P(2) + P'(2)(4 - 2) \; \approx \; 7.
      $$
    \end{itemize}
    Which approximation do you think is a better approximation?  What evidence can you provide to support your claim?
  \end{problem}
    
  \begin{solution} 
    Student {\#2}'s estimate is better.  
    Student {\#1}'s estimate is an underestimate because it assumes 
    the population density grows at a {\em constant} rate for the four years, whereas the sketch in Part (b) 
    suggests that the population density curve is initially concave up.
    The first student's estimate is based upon an assumption that population density
    grows at a constant rate for the first four years.  The second student's estimate is
    based upon an assumption that population density grows at a constant rate only for the
    two years.  Based upon their projection of population density at $t = 2$, they then
    assume that population density grows at a slightly different constant rate during the
    following two years.
%    \includegraphics[width=2.5in]{ws/images/qualitative-eulers-method-turtles-twostep}
  \end{solution}
    
  \item
  \begin{problem}[1in]
    How might you go about obtaining a more accurate estimate for $P(4)$?
    This doesn't need to be something that you carry out by hand, but perhaps 
    some instructions you might give to a computer.
  \end{problem} 
      
  \begin{solution}
    One great way to improve accuracy is to take smaller time steps.
    For instance, we are given the initial condition $P(0) = 2$, from which
    we also know that $P'(0) = 1$.  To estimate $P(0.1)$, we can 
    use linear approximation:
    $$
      P(0.1) \approx P(0) + P'(0) \cdot (0.1 - 0).
    $$
    The number on the right-hand side (call it $y_1$) isn't the exact value of 
    $P(0.1)$, but it's a reasonable estimate.  
    To generate an estimate of $P(0.2)$, we'll use the linear approximation
    \begin{equation}
      \label{second-iterate}
      P(0.2) \approx P(0.1) + P'(0.1) \cdot (0.2 - 0.1)
    \end{equation}
    We already have an estimate for $P(0.1)$, the number we called $y_1$ above.
    The differential equation then gives us an estimate for $P'(0.1)$:
    $$
      P'(0.1) = \frac{5}{8} P(0.1) \left[ 1 - \frac{P(0.1)}{10} \right]
      \; \approx \; \frac{5}{8} y_{1} \left[ 1 - \frac{y_1}{10} \right].
    $$
    Thus, inserting our estimates for $P(0.1)$ and $P'(0.1)$ 
    into Equation~(\ref{second-iterate}) above gives
    \begin{equation}
      \label{second-iterate2}
      P(0.2) \approx y_{1} +  \frac{5}{8} y_{1} \left[ 1 - \frac{y_1}{10} \right] \cdot (0.2 - 0.1).
    \end{equation}
    Let $y_2$ denote the right-hand side of Equation~(\ref{second-iterate2}).
    Repeating this process recursively, we generate a list of estimates
    $$
      y_1 \approx P(0.1), \qquad y_{2} \approx P(0.2), \qquad \dots \qquad
      y_k \approx P(0.1k).
    $$
    After 40 rounds of iteration, we would have $y_{40}$ as our estimate for $P(4)$.      

  \fbox{The name we give to this recursive process is \textbf{Euler's Method}.}

  In the example we just outlined, we refer to $0.1$ as the \textbf{step size}.
  \end{solution}
  
\end{enumerate}



\end{enumerate}


\end{document}






%%%%%%%%%%%%%%%%%%%%%%%%%%%%%%%%%%%%%%%%%%%%%%%%%%%%%%%%%%%%%%%%%%%%%%%%%%%%%%%%%%%%
%%%%%%%%%%%%%%%%%%%%%%%%%%%%%%%%%%%%%%%%%%%%%%%%%%%%%%%%%%%%%%%%%%%%%%%%%%%%%%%%%%%%
%%%%%%%%%%%%%%%%%%%%%%%%%%%%%%%%%%%%%%%%%%%%%%%%%%%%%%%%%%%%%%%%%%%%%%%%%%%%%%%%%%%%
%%%%%%%%%%%%%%%%%%%%%%%%%%%%%%%%%%%%%%%%%%%%%%%%%%%%%%%%%%%%%%%%%%%%%%%%%%%%%%%%%%%%
%%%%%%%%%%%%%%%%%%%%%%%%%%%%%%%%%%%%%%%%%%%%%%%%%%%%%%%%%%%%%%%%%%%%%%%%%%%%%%%%%%%%
%%%%%%%%%%%%%%%%%%%%%%%%%%%%%%%%%%%%%%%%%%%%%%%%%%%%%%%%%%%%%%%%%%%%%%%%%%%%%%%%%%%%
%%%%%%%%%%%%%%%%%%%%%%%%%%%%%%%%%%%%%%%%%%%%%%%%%%%%%%%%%%%%%%%%%%%%%%%%%%%%%%%%%%%%
%%%%%%%%%%%%%%%%%%%%%%%%%%%%%%%%%%%%%%%%%%%%%%%%%%%%%%%%%%%%%%%%%%%%%%%


\item Population growth: Let $P = P(t)$ be the population of Koi in a pond. 

\begin{enumerate}

\item
\begin{problem}
  Suppose you have a population of Koi fish (beautiful orange carp) that, 
  under ideal conditions, grows proportional to itself.  In a Japanese 
  garden Koi are put in a pond. Why do you {\it not} expect the 
  population to grow at a rate proportional to itself in the long run?
\end{problem}

\begin{solution}
  Let $P(t)$ denote population at time $t$.  If the population grew at
  a rate proportional to itself, then $\frac{dP}{dt} = k P$ for some
  positive constant $k$.  From our last problem set, we know that the
  general solution of this differential equation is $P(t) = Ce^{kt}$.
  If you set $t = 0$, you'll find that $C = P(0)$, which is the initial
  population.  Consequently, this model predicts that population grows
  exponentially forever, which is not sustainable in a pond of finite
  size.   
\end{solution}

\item
\begin{problem}
  How might you expect the population to grow in the pond?  Sketch a 
  graph and draw a slope field.  If the pond has a carrying capacity for 
  Koi, what do you think the slope field should look like when the 
  population is below the carrying capacity?  Above the carrying capacity? 
  At the carrying capacity?
\end{problem}

\begin{solution}
  Suppose that $M$ denotes the maximum sustainable population (i.e., the
  carrying capacity).  If $P(t) > M$, then we would expect $P$ to 
  decrease, and if $P(t) < M$ then the environment could sustain more
  Koi (meaning that $P$ can still grow).  Here is what a slope field
  might look like, along with a few representative solution curves.
  \begin{center}
    \includegraphics[width=2.5in]{ws/images/draw-slope-fields-logistic}
  \end{center}
\end{solution}

\item
\begin{problem}
  Use your answers to the previous question to write a differential equation 
  modeling $P$, the population of Koi in a pond with limited space and resources.
\end{problem}

\begin{solution}
  One equation that works is 
  $$
    \frac{dP}{dt} = kP \left( 1 - \frac{P}{M} \right),
  $$
  which is called the {\em logistic equation} and was originally proposed
  by Verhulst in the early 1800s.

  Here's one way to motivate where the logistic equation comes from.  We
  know that the model $\frac{dP}{dt} = kP$ is not feasible, due to limitations
  of resources.  Instead, there ought to be a ``shutoff mechanism" to prevent
  out-of-control growth, and the ``shutoff mechanism" ought to take into
  account the population.  So let's modify $\frac{dP}{dt} = k P$ by
  replacing the relative growth rate $k$ by a population-dependent growth
  rate $k f(P)$.  The differential equation becomes $\frac{dP}{dt} = k P \cdot f(P)$.
  Now let's come up with a reasonable choice for the function $f(P)$.  For
  small populations, exponential growth might actually be sustainable:
  if $P$ is close to zero, we might expect $f(P) \approx 1$.  More 
  specifically, we'll propose a function $f$ for which $f(0) = 1$.
  Next, note that if $P = M$, (the carrying capacity), we would want $f$
  to be zero:  that is, $f(M) = 0$.  Arguably the simplest example of a 
  function $f$ such that $f(0) = 1$ and $f(M) = 0$ is the straight line:
  $f(P) = 1 - P/M$.  This choice of $f$ yields the logistic equation.

  {\em Remark:}  There are other ways to present the logistic equation,
  such as $\frac{dP}{dt} = k P (M - P)$.  There's nothing wrong with writing
  the logistic equation in this way, but the disadvantage is that $k$ will
  have weird units.  Can you figure out what those units would be?
\end{solution}


\end{enumerate}

%%%%%%%%%%%%%%%%%%%%%%%%%%%%%%%%%%%%%%%%%%%%%%%%%%%%%%%%%%%%%%%%
\item 

\note{Students often try to solve partway to write an equation for
$s'(t)$ instead of just writing the equation for $s''(t)$.  It's
worthwhile to note that we have two initial conditions here and that
this has to do with the fact that our differential equation involves a
second derivative.}

\begin{enumerate}
  \item 
\begin{problem}[.7in]
Write a differential equation that models the following situation:

A yellow rubber duck is dropped out of the window of an apartment
building from a height of 80 m.  Let $s = s(t)$ be the height of the
duck above the ground at time $t$.  (Gravity is the acceleration $-9.8 
\textrm{ m/s}^2$.)

\probonly{\textit{Hint:} Write an equation involving a second
derivative.} 
\end{problem}  

\begin{solution}
Due to gravity, the rubber duck has a constant acceleration of $-9.8$
m/s$^2$.  Therefore, $\boxed{s''(t) = -9.8}$.  We know that the rubber
duck starts at 80 feet above the ground.  If we let $t = 0$ be the
time that the rubber duck is dropped, then \fbox{$s(0) = 80$}.  We
also know that the duck is not moving at the beginning, so its initial
velocity is 0.  That is, \fbox{$s'(0) = 0$}.
\end{solution}

\item
\begin{problem}
  How many initial conditions did you include in your solution to part (a)?  Why does it make sense that you'd need this many to give a complete description of the duck's height over time?
\end{problem}
\begin{solution}
  We needed both the initial position $s(0)=80$, and the initial velocity $s'(0)$.
\end{solution}

A differential equation is called `first-order' if it involves only the first derivative of the quantity it models.  It's called second order if it involves only the first and second derivatives, etc..
\end{enumerate}




%%%%%%%%%%%%%%%%%%%%%%%%%%%%%%%%%%%%%%%%%%%%%%%%%%%%%%%%%%%%%%%%%%%%%%%
\item  Consider the differential equation $\frac{dM}{dt} =
    0.1M - 100$.
\begin{enumerate}
\item 
\begin{problem}[1.in]
    Draw the slope field for the differential equation $\frac{dM}{dt} =
    0.01M - 100$.  Sketch several solutions to the differential equation.    
\end{problem}

\begin{solution}
    Here is a sketch of the slope field, along with two representative
    solution curves.
    \includegraphics[width=2.5in]{ws/images/draw-slope-fields-bilbo}
\end{solution}

\item     
\begin{problem}[1.35in]
  Which of the following might be a scenario modeled by this differential equation,
  assuming that $t$ is measured in weeks and $M(t)$ represents population of fish in
  a lake?
  \begin{enumerate}
     \item In any given week, only about 1\% of the fish will survive, and an
     additional 100 fish per week are caught and removed from the lake.
     \item If there were no fishing, then the population would increase by about
     1\% in a given week.  But the rate of fishing increases linearly, and about
     $100t$ fish are removed during Week $t$.
     \item If there were no fishing, then the population would increase by about
     1\% in a given week.  But fish are removed at a constant rate of 100 fish per week.
  \end{enumerate}
\end{problem}

\begin{solution}
  The last option is the correct choice.  Remember that each term in the differential
  equation has to have the same units.  In this case, this means fish per week.  Thus,
  the 100 has units of fish per week---there is a constant harvesting rate of 100 fish
  per week.
\end{solution}


\item     
\begin{problem}[1.35in]
  How do the different solutions you sketched in part (a) relate to the scenario in part (b)?
\end{problem}

\begin{solution}
  If the initial population is more than 1000, then 
  the population will grow over time---the rate of accumulation of fish will outpace
  the continuous withdrawal of fish from the lake.  If the initial population
  in the lake is less than 1000, then the fish population will be drained eventually.
  If the initial population in the lake is exactly 1000, then the
  population will remain constant (because the rate in is always exactly equal to 
  the rate out).
\end{solution}





    
\end{enumerate}

\end{document}